\section{特斯拉 Optimus：一线研发的观察}

\subsection{与 Elon Musk 的工作体验}

杨硕于2023年博士毕业后加入特斯拉Optimus团队，工作约一年半。他对Elon Musk的工作风格有以下观察：

\begin{itemize}
    \item \textbf{高度专注}：开会时基本不开小差、不跑神，非常认真地看演讲者并 follow 整个过程
    \item \textbf{信息收集效率极高}：从一线工程师处直接收集技术进展，每个工程师上来讲几分钟就换下一个，一个半小时可以听二三十个工程师的汇报
    \item \textbf{宏观决策}：最后从宏观角度给项目负责人布置大方向调整，而非纠缠于技术细节
    \item \textbf{对Optimus项目保持耐心}：知道每个细节都很前沿、确定性不高，对整件事情有很好的把控
\end{itemize}

关于"汇报时被开除"的传言，杨硕表示Optimus项目没有碰到过这种情况，他理解这可能发生在一些处于生死关头或需要重大决策的项目中。

\begin{knowledgebox}{Tesla Optimus 的体系化优势}
杨硕认为 Tesla 在机器人可靠性方面最接近波音水平，核心原因是其机器人业务继承了汽车业务线已建立的完整基础设施。如果 Tesla 年内造出5000台机器人，不管部署在哪里，其故障收集、故障分析都有非常完善的已有体系和流程。这种体系化能力是一般创业公司短期内难以复制的。
\end{knowledgebox}

\subsection{大疆与特斯拉的使命感}

杨硕在大疆和特斯拉都深刻感受到了一种\textbf{使命感}驱动的企业文化。两家公司的共同点是：

\begin{itemize}
    \item 始终追求创造全新的产品品类和用户体验
    \item 大疆把无人机和手持相机做成了全新品类
    \item 特斯拉把电动汽车做好做新，Elon在内部不断强调"改变人类未来出行方式"的使命感
    \item 做人形机器人时，Elon 始终站在"整个人类的未来"的角度思考问题
\end{itemize}

杨硕认为，这种使命感是两家公司成功的重要原因——即使商业价值是事后才显现的，但正是使命感驱动了创新产品的诞生。

\subsection{Optimus 的技术进展}

在Tesla内部，杨硕已经看到机器人具备了多项基本能力的迹象，包括：基本的性能表现、简单行走、遥操作叠衣服、自动捡放物体等。他判断可能用不了五年，甚至一两年内，这些就会成为有实际应用价值的产品。

但他也清醒地指出，"做出来"和"能卖出去"之间仍有不小的差距，关键在于场景。例如拿一瓶饮料和拿一个小乐高积木的难度完全不同；如果应用场景是拿放比较大的物体（如亚马逊快递盒），Agility Robotics 的机器人已经演示了六七个小时不停搬运盒子且零失败的表现，已经非常接近商业应用。但如果你期望的是"买一个机器人放在家里，说一句话它就把桌子收了"，这在五到十年内可能还无法实现。

杨硕还指出了一个有趣的观察：所有机器人学家在预测未来时都会说"五年"——但实际发展往往更长。不过他自己在Tesla内部看到的进展让他相信，某些特定场景下的机器人产品可能真的只需要一两年。

\begin{knowledgebox}{"五年预测"的心理惯性}
杨硕指出一个有趣的行业现象：几乎所有机器人学家被问及"机器人何时走进家庭"时，都会回答"五年"。这反映了一种心理惯性——五年既足够远以至于不用负责，又足够近以至于让人兴奋。杨硕虽然也给出了类似的时间框架，但他基于在Tesla内部的实际观察，对特定垂直场景持更乐观的态度，同时对通用家庭助手场景保持谨慎。
\end{knowledgebox}

\subsection{本章小结}

特斯拉的经历让杨硕亲眼见证了大规模工业体系如何赋能前沿机器人研发。Elon的高效信息收集与宏观决策风格、Tesla从汽车业务继承的体系化基础设施、以及使命感驱动的企业文化，都是Optimus项目独特竞争力的来源。


